\documentclass[10pt,a4paper]{article}
\usepackage[margin=0.6in]{geometry}
\usepackage{amsmath,amssymb,amsfonts,stmaryrd}
\usepackage{xcolor}
\usepackage{enumitem}
\usepackage{tcolorbox}
\usepackage{multicol}
\usepackage{tikz}
\usetikzlibrary{arrows.meta, positioning}
\usepackage{listings}
\usepackage{hyperref}

\definecolor{isabelleblue}{RGB}{0, 70, 140}
\definecolor{isabellekw}{RGB}{120, 0, 120}
\definecolor{boxbg}{RGB}{248, 249, 250}

\tcbset{
  colback=boxbg,
  colframe=isabelleblue,
  fonttitle=\bfseries\large,
  boxrule=0.8pt,
  arc=3pt,
  left=8pt,
  right=8pt,
  top=6pt,
  bottom=6pt
}

\title{\textbf{\huge Mathematical Induction for Natural Numbers}\\[0.3em]
\large Lecture Notes}
\author{\textbf{Course:} Computer-Assisted Theorem Proving (7CCSACTP)}
\date{}

\begin{document}

\maketitle
\vspace{-1.5em}

\begin{tcolorbox}[title=Context \& Induction Intuition]
\begin{multicols}{2}
\textbf{$\rightarrow$ Previous Videos:}
\begin{itemize}[leftmargin=1.2em, itemsep=2pt]
    \item Basic elements of interface (terms, theorems)
    \item Basic proof methods (FW, BW, \texttt{subst})
\end{itemize}

\textbf{$\rightarrow$ \underline{This Video}:}
\begin{itemize}[leftmargin=1.2em, itemsep=2pt]
    \item Mathematical induction for natural numbers (\texttt{nat})
\end{itemize}

\columnbreak

\textbf{Intuition for Mathematical Induction:}
\begin{itemize}[leftmargin=1.2em, itemsep=2pt]
    \item To prove a property $P$ holds for all $n \in \mathbb{N}$ (e.g. $?x \le ?x + 3$), we prove two cases:
    \begin{itemize}[leftmargin=1em, label=$\bullet$]
        \item $\textbf{goal}_1 \textbf{ (Base Case):}$ $P(0)$ holds.
        \item $\textbf{goal}_2 \textbf{ (Step Case):}$ Assuming $P(n)$ holds (Induction Hypothesis, \textsf{I.H.}), $P(n+1)$ holds.
    \end{itemize}
    \item Algorithmic view of induction:
    \begin{center}
    $P(0) \xrightarrow{\textbf{goal}_1} P(1) \xrightarrow{\textbf{goal}_2} P(2) \xrightarrow{\textbf{goal}_2} \dots \xrightarrow{\textbf{goal}_2} P(n)$
    \end{center}
    \item \textbf{Isabelle Note:} In Isabelle, $n+1$ is represented as \texttt{Suc n} for natural numbers.
\end{itemize}
\end{multicols}
\end{tcolorbox}

\vspace{0.3em}

\begin{tcolorbox}[title={Detailed Example: Commutativity of Addition ($?x + ?y = ?y + ?x$)}]
\textbf{Proof Strategy:} Induction on $y$.

\begin{multicols}{2}
\textbf{Base Case ($y = 0$):}
\begin{itemize}[leftmargin=1.2em, itemsep=2pt]
    \item $\textbf{goal}_1::$ $x + 0 = 0 + x$
    \item \textit{Lemmas:} \texttt{lem1:} $x + 0 = x$, \quad \texttt{lem2:} $0 + x = x$
    \item \textbf{Proof steps:}
    \begin{enumerate}[leftmargin=1.2em, itemsep=1pt]
        \item $x + 0 = x$ \quad \textit{(by subst \texttt{lem1})}
        \item $\dots = 0 + x$ \quad \textit{(by subst \texttt{lem2})}
    \end{enumerate}
    \item \textsf{finally} $x + 0 = 0 + x \quad \checkmark$
\end{itemize}

\columnbreak

\textbf{Step Case ($y \to \text{Suc } y$):}
\begin{itemize}[leftmargin=1.2em, itemsep=2pt]
    \item \textbf{Assume I.H.:} $x + y = y + x$
    \item $\textbf{goal}_2::$ $x + (\text{Suc } y) = (\text{Suc } y) + x$
    \item \textit{Lemmas:}
    \begin{itemize}[leftmargin=1em, label=$\bullet$]
        \item \texttt{lem3:} $?a + (\text{Suc } ?b) = \text{Suc } (?a + ?b)$
        \item \texttt{lem4:} $(\text{Suc } ?a) + ?b = \text{Suc } (?a + ?b)$
    \end{itemize}
    \item \textbf{Proof steps:}
    \begin{enumerate}[leftmargin=1.2em, itemsep=1pt]
        \item $x + (\text{Suc } y) = \text{Suc } (x + y)$ \quad \textit{(by subst \texttt{lem3})}
        \item $\dots = \text{Suc } (y + x)$ \quad \textit{(by subst \textsf{I.H.})}
        \item $\dots = (\text{Suc } y) + x$ \quad \textit{(by subst \texttt{lem4})}
    \end{enumerate}
    \item \textsf{finally} $x + (\text{Suc } y) = (\text{Suc } y) + x \quad \checkmark$
\end{itemize}
\end{multicols}
\end{tcolorbox}

\vspace{0.3em}

\begin{tcolorbox}[title=Isabelle Syntax Reminder]
\begin{itemize}[leftmargin=1.2em, itemsep=2pt]
    \item In Isabelle: \texttt{Suc x} $\approx x + 1$ for terms of type \texttt{nat}.
    \item Induction rule applied via: \texttt{apply (induct y)} or \texttt{proof (induct y)}.
\end{itemize}
\end{tcolorbox}

\end{document}
